\chapter{Modules over a Principal Ideal Domain}

\begin{center}\itshape
How many abelian groups of order $360$ are there?\\
And why does every square matrix satisfy its own characteristic polynomial?\\
One theorem answers both: the structure of finitely generated modules over a principal ideal domain.
\end{center}

How many different abelian groups have $360$ elements? The number factors as
\[
360=2^3\cdot3^2\cdot5.
\]
There are three ways to split the $2^3$ part into cyclic prime-power factors, two ways for $3^2$, and one for $5$. That gives $3\cdot2=6$ groups up to isomorphism. Why should counting partitions of exponents classify groups?

There is another question that seems unrelated. Every square matrix $A$ satisfies its own characteristic polynomial:
\[
\det(XI-A)\big|_{X=A}=0.
\]
Why can an expression involving a determinant turn into the zero matrix when we substitute the matrix for $X$?

Both questions become easier in the language of \emph{modules}. An abelian group is a module over $\Z$: multiplying by an integer means repeated addition. A vector space with a linear transformation $A$ is a module over $K[X]$: the polynomial $P$ acts as $P(A)$. The underlying idea is the same, even though one case starts with groups and the other with matrices.

Over a principal ideal domain, finitely generated modules have a precise structure. The Smith normal form is the algorithmic side of that structure; the classification of finite abelian groups is one consequence, and rational and Jordan canonical forms are others.

We will prove the structure statement rather than treat it as a black box. Watch how division with remainder turns a matrix problem into a list of invariant factors. Those factors will continue to matter when we study field extensions.

\section{Modules}

\begin{definition}
An \emph{$R$-module} is an abelian group $(M,+)$ with a map $R\times M\to M$, $(r,x)\mapsto rx$, such that $r(x+y)=rx+ry$, $(r+s)x=rx+sx$, $(rs)x=r(sx)$ and $1x=x$. Submodules, quotients $M/N$, homomorphisms (additive maps with $f(rx)=rf(x)$), kernels and images are defined as for vector spaces, and the first isomorphism theorem $M/\ker f\cong f(M)$ holds with the same proof (Theorem~\ref{thm:iso1}).
\end{definition}

\begin{example}\label{ex:modules}
\begin{enumerate}[label=(\arabic*)]
\item Over a field, modules are vector spaces.
\item Over $\Z$, modules are abelian groups: $nx$ must be $x+\dots+x$, so the module structure is forced.
\item The ideals of $R$ are its submodules; $R/I$ is a module.
\item $R^n$, with componentwise operations. A module is \emph{free} if it has a \emph{basis}, a family $(e_i)$ such that every element is uniquely $\sum r_ie_i$ (finitely many $r_i\ne0$); a free module with a basis of $n$ elements is isomorphic to $R^n$.
\item If $V$ is a vector space over $K$ and $A:V\to V$ is linear, $V$ becomes a $K[X]$-module $V_A$ by $P\cdot v=P(A)v$. Submodules are the $A$-stable subspaces.
\end{enumerate}
\end{example}

\begin{think}
Why should a matrix be ``the same thing'' as a $K[X]$-module? If $V_A\cong V_B$ as $K[X]$-modules, by a bijection $S$, what does $S$ do to $A$ and $B$?
\end{think}

A module is \emph{finitely generated} if it equals $Rx_1+\dots+Rx_n$ for some $x_i$, equivalently if there is a surjective homomorphism $R^n\to M$, $(r_i)\mapsto\sum r_ix_i$. If $R$ is a domain, the number of elements of a basis of a free module is well defined: an isomorphism $R^s\to R^t$ and its inverse are given by matrices $B$ ($t\times s$) and $C$ ($s\times t$) with $CB=I_s$, $BC=I_t$; these are also inverse linear maps between $K^s$ and $K^t$, $K=\operatorname{Frac}R$, so $s=t$. This number is the \emph{rank}.

\subsection*{An application: the norm of an ideal is its index}

Volume~13, Chapter~1 defined the norm $N(I)$ of a non-zero ideal $I$ of the ring of integers $\mathcal O$ of a quadratic field by $I\bar I=(N(I))$, and announced that $N(I)$ is the number of classes of $\mathcal O$ modulo $I$. Module language proves it in a few lines, using the results of that chapter (cancellation, ``to contain is to divide'', and factorisation into maximal ideals).

\begin{theorem}\label{thm:normindex}
For every non-zero ideal $I$ of $\mathcal O$, $[\mathcal O:I]=N(I)$.
\end{theorem}
\begin{proof}
\emph{Maximal ideals.} Let $P$ be maximal. It contains the positive integer $N(P)$, hence, being prime (Volume~13, Chapter~1: if $xy\in P$ then $x\in P$ or $y\in P$), one of its prime factors $p$. Then $p\mathcal O\subseteq P\subsetneq\mathcal O$, and $[\mathcal O:p\mathcal O]=p^2$ since $\mathcal O\cong\Z^2$ as a group. If $P=p\mathcal O$, then $[\mathcal O:P]=p^2=N(p\mathcal O)=N(P)$. Otherwise $p\mathcal O\subsetneq P\subsetneq\mathcal O$, so $[\mathcal O:P]$ is a divisor of $p^2$ other than $1$ and $p^2$: it is $p$. Also $P\supseteq p\mathcal O$ gives $p\mathcal O=PL$ with $L\ne\mathcal O$ (else $P=p\mathcal O$), so $N(P)N(L)=p^2$ with both factors $>1$: $N(P)=p$.

\emph{One more factor.} Let $I$ be an ideal and $P$ maximal; we show $[I:IP]=[\mathcal O:P]$. The quotient $I/IP$ is an $\mathcal O$-module killed by $P$, so it is a vector space over the field $k=\mathcal O/P$ (Proposition~\ref{prop:primemax}), and its subspaces are the $L/IP$ with $L$ an ideal, $IP\subseteq L\subseteq I$. For such $L$, write $L=IL'$ (to contain is to divide); $IP\subseteq IL'$ gives $IP=IL'L''$, so $P=L'L''$ by cancellation, $L'\supseteq P$, and $L'\in\{P,\mathcal O\}$. So the only subspaces are $0$ and $I/IP$, which is non-zero ($IP=I$ would give $P=\mathcal O$ by cancellation). A non-zero vector space with no other subspace has dimension $1$, so $|I/IP|=|k|$.

\emph{Conclusion.} Write $I=P_1\cdots P_r$ (Dedekind's theorem). By induction on $r$, using $[\mathcal O:IP]=[\mathcal O:I]\,[I:IP]$ and the multiplicativity of $N$, $[\mathcal O:I]=N(P_1)\cdots N(P_r)=N(I)$.
\end{proof}

A module can be described by generators and the relations among them. Writing those relations as a matrix lets us simplify the description using allowed row and column operations. The Smith normal form records what survives.

\section{The Smith normal form}\label{sec:smith}

From now on $R$ is a PID. A square matrix over $R$ is invertible over $R$ iff its determinant is a unit (the adjugate formula $B\operatorname{adj}B=(\det B)I$ holds over any commutative ring); these matrices form the group $GL_n(R)$. Row operations (swapping two rows, adding $r$ times one row to another, multiplying a row by a unit) are multiplications on the left by matrices of $GL_m(R)$, and column operations are multiplications on the right.

\begin{theorem}[Smith normal form]\label{thm:smith}
Let $A$ be an $m\times n$ matrix over a PID $R$. There are $P\in GL_m(R)$ and $Q\in GL_n(R)$ such that
\[
\begin{gathered}PAQ=\begin{pmatrix}D&0\\0&0\end{pmatrix},\qq D=\operatorname{diag}(d_1,\dots,d_r),\\ d_i\ne0,\quad d_1\mid d_2\mid\dots\mid d_r.\end{gathered}
\]
The number $r$ and the $d_i$ (up to units) are determined by $A$: if $\Delta_k(A)$ denotes a gcd of the $k\times k$ minors of $A$, then $d_1\cdots d_k=\Delta_k(A)$ up to a unit.
\end{theorem}

\begin{proof}
\emph{Existence.} For $a\ne0$ let $\lambda(a)$ be the number of prime factors of $a$, with multiplicity ($R$ is a UFD by Theorem~\ref{thm:edpidufd}). If $A=0$ there is nothing to do. Otherwise move an entry with smallest $\lambda$ to position $(1,1)$.

\emph{Gcd step.} Suppose $a=a_{11}$ does not divide some $b=a_{i1}$. Let $(a,b)=(d)$, $d=ua+vb$, $a=da'$, $b=db'$, so that $ua'+vb'=1$. The matrix $\begin{psmallmatrix}u&v\\-b'&a'\end{psmallmatrix}$ has determinant $1$ and sends $\binom ab$ to $\binom d0$; acting with it on rows $1$ and $i$ replaces $a_{11}$ by $d$, a proper divisor of $a$ (not an associate, since $a\nmid b$ and $d\mid b$), so $\lambda(d)<\lambda(a)$. The same works on columns if $a_{11}$ does not divide some $a_{1j}$.

\emph{Clearing.} If $a_{11}$ divides every entry of the first row and column, subtracting multiples of the first row and column makes them zero except $a_{11}$. If now $a_{11}$ does not divide some entry $a_{ij}$ ($i,j\ge2$), add row $i$ to row $1$ and do a gcd step.

Each gcd step lowers $\lambda(a_{11})$, so after finitely many steps $A$ becomes $\begin{psmallmatrix}d_1&0\\0&A'\end{psmallmatrix}$ with $d_1$ dividing every entry of $A'$. Apply the same procedure to $A'$; its operations do not touch the first row and column, and they replace entries by $R$-linear combinations of entries, so all later entries, in particular $d_2,d_3,\dots$, remain multiples of $d_1$. By induction we reach the form claimed, with $d_1\mid d_2\mid\cdots$.

\emph{Uniqueness.} The rows of $PA$ are $R$-linear combinations of the rows of $A$; since a determinant is multilinear and alternating in its rows, every $k\times k$ minor of $PA$ is an $R$-linear combination of $k\times k$ minors of $A$, so $\Delta_k(A)$ divides $\Delta_k(PA)$. The same holds for columns and $AQ$. As $P,Q$ are invertible, $\Delta_k(A)$ and $\Delta_k(PAQ)$ divide each other. For the diagonal matrix, a non-zero $k\times k$ minor is a product $d_{i_1}\cdots d_{i_k}$ with $i_1<\dots<i_k\le r$, which is a multiple of $d_1\cdots d_k$ because $d_j\mid d_{i_j}$; so $\Delta_k=d_1\cdots d_k$ for $k\le r$ and $\Delta_k=0$ for $k>r$. Hence $r$ and $d_k=\Delta_k/\Delta_{k-1}$ are determined up to units.
\end{proof}

\paragraph{Why this works.}
Think of a presentation matrix as a list of relations. We may replace one relation by a sum of relations, or replace the generators by another basis, without changing the module being presented. The Smith algorithm repeatedly uses division with remainder to put the smallest non-zero relation into the top-left position and make it divide the others. The final diagonal entries are not just convenient numbers: they describe the cyclic pieces of the quotient module. This is why the same algorithm works for an abelian group presented over $\Z$ and for a linear operator presented over $K[X]$.

Over a Euclidean ring such as $\Z$ or $K[X]$, gcd steps can be replaced by Euclidean divisions, and the proof is an algorithm, carried out in Figure~\ref{fig:smith} .

\begin{figure}[ht]
\centering
\fig{ch07-smith}
\caption{The Smith normal form of an integer matrix, by row operations ($R_i$) and column operations ($C_j$). The invariants can be checked with minors: $\Delta_1=\gcd$ of the entries $=2$, $\Delta_3=|\det A|=144=2\cdot6\cdot12$.}
\label{fig:smith}
\end{figure}

The diagonal entries of a Smith form do more than simplify a calculation. They tell us what pieces the whole module is made from. Over a principal ideal domain, those pieces are a free part and cyclic torsion factors.

\section{The structure theorem}

\begin{lemma}\label{lem:subfree}
Over a PID, every submodule of $R^n$ is free of rank at most $n$.
\end{lemma}
\begin{proof}
Induction on $n$; for $n=0$ it is clear. Let $K\subseteq R^n$ be a submodule. The last coordinates of the elements of $K$ form an ideal $(a)$ of $R$. The elements of $K$ with last coordinate $0$ form a submodule $K'$ of $R^{n-1}$, free of rank $\le n-1$ by induction. If $a=0$, $K=K'$. Otherwise pick $k\in K$ with last coordinate $a$. Any $x\in K$ has last coordinate $ca$, and $x-ck\in K'$; and $K'\cap Rk=0$, because $rk$ has last coordinate $ra$, which is $0$ only if $r=0$. So $K=K'\oplus Rk$, and a basis of $K'$ together with $k$ is a basis of $K$.
\end{proof}

\begin{theorem}[structure of finitely generated modules over a PID]\label{thm:structure}
Let $M$ be a finitely generated module over a PID $R$. Then
\[
M\cong R^s\oplus R/(d_1)\oplus\dots\oplus R/(d_r),
\]
with $d_1\mid d_2\mid\dots\mid d_r$ non-zero non-units. The integer $s$ (the \emph{rank}) and the \emph{invariant factors} $d_i$, up to units, are determined by $M$.
\end{theorem}

\begin{proof}
\emph{Existence.} Choose generators $x_1,\dots,x_n$ and the surjection $\varphi:R^n\to M$, $e_i\mapsto x_i$. By Lemma~\ref{lem:subfree}, $\ker\varphi$ has a basis $k_1,\dots,k_m$, $m\le n$; let $A$ be the $n\times m$ matrix with columns $k_j$. Take $P,Q$ as in Theorem~\ref{thm:smith}. The columns of $AQ$ are again a basis of $\ker\varphi$, and $AQ=P^{-1}(PAQ)$. Let $f_i$ be the $i$-th column of $P^{-1}$; the $f_i$ form a basis of $R^n$, and the $j$-th column of $AQ$ is $d_jf_j$ for $j\le r$, and $0$ for $j>r$. Hence $\ker\varphi=Rd_1f_1\oplus\dots\oplus Rd_rf_r$, and
\[
\begin{aligned}M\cong R^n/\ker\varphi&\cong R/(d_1)\oplus\dots\oplus R/(d_r)\eqbr\oplus R^{n-r}.\end{aligned}
\]
Drop the $d_i$ that are units, since then $R/(d_i)=0$.

\emph{Uniqueness.} The \emph{torsion} submodule $T=\{x\in M:ax=0\text{ for some }a\ne0\}$ is $\bigoplus R/(d_i)$ (the elements of $R^s$ killed by some $a\ne0$ are $0$, $R$ being a domain), and $M/T\cong R^s$; isomorphic modules have isomorphic torsion submodules and quotients, so $s$ is determined. For $T$, fix a prime $p$ of $R$ and $k\ge0$. Then $p^kT/p^{k+1}T$ is killed by $p$, so it is a vector space over the field $R/(p)$, and it is the direct sum of the spaces $p^kN/p^{k+1}N$ for $N=R/(d_i)$. If $d=p^ve$ with $p\nmid e$ and $N=R/(d)$, then $p^kN=(p^k,d)/(d)=(p^{\min(k,v)})/(d)$; so for $k\ge v$, $p^kN=p^{k+1}N$, and for $k<v$, $p^kN/p^{k+1}N\cong(p^k)/(p^{k+1})\cong R/(p)$ (multiplication by $p^k$). Hence
\[
\dim_{R/(p)}\bigl(p^kT/p^{k+1}T\bigr)=\#\{i:v_p(d_i)>k\}.
\]
These numbers, for all $k$, determine the non-zero exponents $v_p(d_i)$ as a multiset. Since $v_p(d_1)\le\dots\le v_p(d_r)$, they determine each $v_p(d_i)$ once $r$ is known; and $r$ is the largest of the numbers $\dim(T/pT)$, attained at any prime $p\mid d_1$ (which divides every $d_i$). So each $d_i$ is determined up to a unit.
\end{proof}

\paragraph{Why this works.}
Existence and uniqueness require different ideas. Existence comes from a presentation $R^n\twoheadrightarrow M$ and the Smith form of its relation matrix. Uniqueness comes from quantities we can recover from $M$ itself: its free rank and, for each prime element $p$, the dimensions of the quotients $p^kT/p^{k+1}T$. Each such dimension counts the cyclic factors that still have at least one more power of $p$ left. For example, in $\Z/8\Z\oplus\Z/4\Z$, the dimensions over $\mathbb F_2$ for $k=0,1,2$ are $2,2,1$. We can read the exponents $3$ and $2$ back from that sequence.

For $R=\Z$, $\Z/d\Z\cong\bigoplus_p\Z/p^{v_p(d)}\Z$ by the Chinese remainder theorem, which gives two equivalent forms of the classification.

\begin{corollary}[finitely generated abelian groups]\label{cor:abelian}
Every finitely generated abelian group is isomorphic to $\Z^s\oplus\Z/d_1\Z\oplus\dots\oplus\Z/d_r\Z$ with $1<d_1\mid\dots\mid d_r$, and also to $\Z^s\oplus\bigoplus_j\Z/p_j^{a_j}\Z$ with prime powers $p_j^{a_j}$; $s$, the $d_i$, and the multiset of the $p_j^{a_j}$ are unique. In particular, if $n=\prod p^{a_p}$, the number of abelian groups of order $n$ is $\prod_p\mathrm{p}(a_p)$, where $\mathrm p(a)$ is the number of partitions of $a$.
\end{corollary}
\begin{proof}
The first form is Theorem~\ref{thm:structure}. Splitting each $\Z/d_i\Z$ by the Chinese remainder theorem gives the second. Conversely, from the multiset of prime powers one recovers the $d_i$: $d_r$ is the product, over $p$, of the largest power of $p$ in the multiset, $d_{r-1}$ the product of the next largest ones, and so on; so the second form is unique too. A group of order $n$ is finite, so $s=0$, and its $p$-part is $\bigoplus_j\Z/p^{a_j}\Z$ with $\sum_ja_j=a_p$: a partition of $a_p$ for each $p$, chosen independently.
\end{proof}

For $360=2^3\cdot3^2\cdot5$ the answer to the opening question is $\mathrm p(3)\,\mathrm p(2)\,\mathrm p(1)=3\cdot2\cdot1=6$ (Exercise~\ref{ex:ab360} lists them).

\begin{corollary}\label{cor:index}
If $A\in M_n(\Z)$ has $\det A\ne0$, the subgroup $A\Z^n$ has index $|\det A|$ in $\Z^n$, and $\Z^n/A\Z^n\cong\bigoplus\Z/d_i\Z$ with the $d_i$ of the Smith form.
\end{corollary}
\begin{proof}
With $PAQ=D$: $A\Z^n=AQ\Z^n$, and $P$ is an automorphism of $\Z^n$ mapping $AQ\Z^n$ onto $D\Z^n$. So $\Z^n/A\Z^n\cong\Z^n/D\Z^n=\bigoplus\Z/d_i\Z$, with $|d_1\cdots d_n|=|\det D|=|\det A|$ elements, since $\det P,\det Q=\pm1$.
\end{proof}

An abelian group is a module over $\Z$. Now let a polynomial act on a vector by evaluating it at a linear transformation. The same structure theorem that classified finite abelian groups can then describe matrices.

\section{Linear algebra over \texorpdfstring{$K[X]$}{K[X]}}

Let $K$ be a field, $A\in M_n(K)$, and $V_A=K^n$ with $P\cdot v=P(A)v$ (Example~\ref{ex:modules}(5)). Two matrices $A,B$ are similar ($B=SAS^{-1}$) iff $V_A\cong V_B$ as $K[X]$-modules: a $K[X]$-isomorphism is a linear bijection $S$ with $S(Av)=B(Sv)$, i.e.\ $SA=BS$. So classifying matrices up to similarity is classifying these modules.

\begin{proposition}\label{prop:presentation}
$V_A\cong K[X]^n/(XI-A)K[X]^n$. Consequently the Smith form of $XI-A$ over $K[X]$ is $\operatorname{diag}(1,\dots,1,d_1,\dots,d_r)$ with $d_i$ monic of positive degree, $d_1\mid\dots\mid d_r$, $d_1\cdots d_r=\chi_A=\det(XI-A)$, and
\[
V_A\cong K[X]/(d_1)\oplus\dots\oplus K[X]/(d_r).
\]
The $d_i$ are the \emph{invariant factors} of $A$.
\end{proposition}
\begin{proof}
Let $\varphi:K[X]^n\to V_A$, $\sum_jF_je_j\mapsto\sum_jF_j(A)e_j$; it is $K[X]$-linear and onto (constant vectors suffice). The $j$-th column of $XI-A$ is $Xe_j-\sum_ia_{ij}e_i$, which $\varphi$ sends to $Ae_j-\sum_ia_{ij}e_i=0$. Conversely, modulo the image of $XI-A$ we have $Xe_j\equiv\sum_ia_{ij}e_i$, a constant vector; hence $X$ times a constant vector is congruent to a constant vector, and by induction on the degree every $F\in K[X]^n$ is congruent to a constant vector $c$. If $\varphi(F)=0$ then $\varphi(c)=c=0$, so $F$ is in the image. Thus $\ker\varphi=(XI-A)K[X]^n$.

$\det(XI-A)=\chi_A\ne0$, so the Smith form has $n$ non-zero diagonal entries; make them monic (multiply by units of $K[X]$, the non-zero constants). Their product is $\det P\det Q\,\chi_A$ with $\det P,\det Q\in K^*$, and both sides are monic, so it equals $\chi_A$. As in the proof of Theorem~\ref{thm:structure}, $V_A\cong\bigoplus K[X]/(d_i)$, the factors $d_i=1$ contributing nothing.
\end{proof}

For a monic $d=X^m+c_{m-1}X^{m-1}+\dots+c_0$, the module $K[X]/(d)$ has the basis $1,x,\dots,x^{m-1}$ ($x$ the class of $X$; use Euclidean division), and multiplication by $x$ has the \emph{companion matrix}
\[
C(d)=\begin{pmatrix}0&&&-c_0\\1&\ddots&&-c_1\\&\ddots&0&\vdots\\&&1&-c_{m-1}\end{pmatrix}.
\]
Its characteristic polynomial is $d$: expanding $\det(XI-C(d))$ along the first row gives $X\det(XI-C(d'))+c_0$ with $d'=X^{m-1}+c_{m-1}X^{m-2}+\dots+c_1$ (the other minor is triangular with $-1$ on its diagonal, determinant $(-1)^{m-1}$, and the sign of position $(1,m)$ is $(-1)^{m+1}$), and induction gives $Xd'+c_0=d$. Also $g(C(d))=0$ iff $g\cdot1=0$ in $K[X]/(d)$ iff $d\mid g$.

\begin{theorem}[rational canonical form; Frobenius]\label{thm:frobenius}
Every $A\in M_n(K)$ is similar to the block diagonal matrix $\operatorname{diag}(C(d_1),\dots,C(d_r))$ of its invariant factors. Two matrices are similar iff they have the same invariant factors. Moreover the minimal polynomial is $\mu_A=d_r$, and $\chi_A=d_1\cdots d_r$.
\end{theorem}
\begin{proof}
In the basis of $V_A\cong\bigoplus K[X]/(d_i)$ made of the bases $1,x,\dots$ of each summand, $A$ (multiplication by $X$) has the stated matrix. Similar matrices have isomorphic modules, hence the same invariant factors by Theorem~\ref{thm:structure}; conversely, same invariant factors give the same block matrix. A polynomial $g$ kills $\bigoplus K[X]/(d_i)$ iff $d_i\mid g$ for all $i$, iff $d_r\mid g$; so $\mu_A=d_r$. The product formula is Proposition~\ref{prop:presentation}.
\end{proof}

\begin{corollary}[Cayley--Hamilton]\label{cor:cayleyhamilton}
$\chi_A(A)=0$. Moreover $\mu_A$ divides $\chi_A$, and they have the same irreducible factors.
\end{corollary}
\begin{proof}
$\mu_A=d_r$ divides $d_1\cdots d_r=\chi_A$, so $\chi_A(A)=0$. An irreducible factor of $\chi_A$ divides some $d_i$, which divides $d_r=\mu_A$.
\end{proof}

When $\chi_A$ splits into linear factors (always if $K=\C$), the Chinese remainder theorem refines the decomposition. Writing $d_i=\prod_\lambda(X-\lambda)^{e_{i,\lambda}}$, Theorem~\ref{thm:crtring} gives $K[X]/(d_i)\cong\bigoplus_\lambda K[X]/\bigl((X-\lambda)^{e_{i,\lambda}}\bigr)$. In $K[X]/((X-\lambda)^e)$, with $y=x-\lambda$, the basis $b_1=y^{e-1},b_2=y^{e-2},\dots,b_e=1$ satisfies $xb_1=\lambda b_1$ and $xb_k=\lambda b_k+b_{k-1}$: multiplication by $x$ is the \emph{Jordan block} $J_e(\lambda)$, with $\lambda$ on the diagonal and $1$ just above it.

\begin{theorem}[Jordan normal form]\label{thm:jordan}
If $\chi_A$ splits over $K$, then $A$ is similar to a block diagonal matrix of Jordan blocks $J_{e}(\lambda)$, unique up to the order of the blocks. The blocks for $\lambda$ have sizes $e_{1,\lambda}\le\dots\le e_{r,\lambda}$ (the non-zero ones), the exponents of $X-\lambda$ in the invariant factors.
\end{theorem}
\begin{proof}
Existence is the computation above. The blocks of a Jordan form of $A$ determine a decomposition $V_A\cong\bigoplus K[X]/((X-\lambda)^e)$, and from such a decomposition the invariant factors are recovered as in the proof of Corollary~\ref{cor:abelian} ($d_r$ is the product of the largest power of each $X-\lambda$, and so on). Two Jordan forms of $A$ therefore have the same multiset of $(X-\lambda)^e$, i.e.\ the same blocks.
\end{proof}

\begin{beauty}
The same theorem classifies abelian groups and matrices, because $\Z$ and $K[X]$ are both PIDs: this is the kind of unification that made abstract algebra worth inventing. It reappears often: the structure of $(\Z/q\Z)^*$ (Exercise~\ref{ex:unitsmodn}) is what Dirichlet characters are built on (Chapter~11, then Volume~6); Jordan forms solve linear recurrences and differential equations (Volume~1) and describe the long-run behaviour of Markov chains (Volume~4); and the Smith form of boundary matrices computes the \emph{homology} of a space, its holes and twists, which a numerical calculation uses to see that the projective plane has a hole of ``order $2$''.
\end{beauty}

\section{Exercises}

\begin{exercise}[$\star$]\label{ex:ab360}
List the six abelian groups of order $360$, in both forms of Corollary~\ref{cor:abelian}.
\end{exercise}

\begin{exercise}[$\star$]\label{ex:smithcalc}
Identify the abelian group with generators $x,y$ and relations $4x+6y=0$, $6x+4y=0$, i.e.\ $\Z^2/A\Z^2$ with $A=\begin{psmallmatrix}4&6\\6&4\end{psmallmatrix}$.
\end{exercise}

\begin{exercise}[$\star\star$]\label{ex:unitsmodn}
\begin{enumerate}[label=(\alph*)]
\item Let $k\ge3$. Show that $5^{2^j}=1+2^{j+2}t_j$ with $t_j$ odd, that $5$ has order $2^{k-2}$ modulo $2^k$, and that $(\Z/2^k\Z)^*\cong\Z/2\Z\times\Z/2^{k-2}\Z$, generated by $-1$ and $5$.
\item Find the invariant factors of $(\Z/360\Z)^*$, and the smallest $e\ge1$ with $a^e\equiv1\pmod{360}$ for all $a$ prime to $360$.
\end{enumerate}
\end{exercise}

\begin{exercise}[$\star\star$]\label{ex:similarQ}
Let $K\subseteq L$ be fields and $A,B\in M_n(K)$. Show that if $A$ and $B$ are similar over $L$, they are similar over $K$. (For instance, two rational matrices similar over $\C$ are similar over $\Q$.)
\end{exercise}

\begin{exercise}[$\star\star$]\label{ex:nilpotent}
Show that the similarity classes of nilpotent $n\times n$ matrices over a field correspond to the partitions of $n$, and that the partition of $N$ can be read from the numbers $\dim\ker N^k$. List the classes for $n=4$.
\end{exercise}

\begin{exercise}[$\star\star$]\label{ex:diagonalisable}
Show that $A\in M_n(K)$ is diagonalisable over $K$ iff $\mu_A$ is a product of distinct factors $X-\lambda$, $\lambda\in K$. Deduce that a matrix with $A^2=A$ is diagonalisable, and that a complex matrix with $A^k=I$ for some $k\ge1$ is diagonalisable.
\end{exercise}

\begin{exercise}[$\star\star$]\label{ex:cyclicvector}
Show that the following are equivalent: (i) $\mu_A=\chi_A$; (ii) there is a vector $v$ such that $v,Av,\dots,A^{n-1}v$ is a basis; (iii) $A$ is similar to a companion matrix.
\end{exercise}

\begin{exercise}[$\star\star\star$]\label{ex:sublattices}
Show that every subgroup of index $n$ of $\Z^2$ can be written uniquely as $\Z(a,0)+\Z(b,c)$ with $a,c\ge1$, $ac=n$ and $0\le b<a$. Deduce that $\Z^2$ has exactly $\sigma(n)=\sum_{a\mid n}a$ subgroups of index $n$.
\end{exercise}

\begin{exercise}[$\star\star\star$]\label{pb:kronecker}
\emph{Kronecker's theorem.}
\begin{enumerate}[label=(\alph*)]
\item Let $P\in\Z[X]$ be monic with $P(0)\ne0$, and suppose that all its complex roots have modulus at most $1$. Show
that they are all roots of unity, and that if $P$ is irreducible, $P$ is a cyclotomic polynomial.
\item Let $P\in\Z[X]$ be monic, with all its roots real and in $[-2,2]$. Show that its roots are of the form $2\cos(2\pi r)$
with $r\in\Q$.
\item The number $z=\frac{3+4i}5$ has modulus $1$, and so does its only conjugate $\bar z$. Show that it is not a root of
unity. Which hypothesis of (a) fails?
\end{enumerate}
\end{exercise}

\section{Hints}

\begin{hint}{ex:ab360}
Partitions of $3$ for the $2$-part, of $2$ for the $3$-part. To pass to invariant factors, multiply the largest powers of each prime together, then the next largest.
\end{hint}
\begin{hint}{ex:smithcalc}
$\Delta_1$ is the gcd of the entries, $\Delta_2=|\det A|$.
\end{hint}
\begin{hint}{ex:unitsmodn}
(a) Square $1+2^{j+2}t_j$. Then $-1\notin\langle5\rangle$ because powers of $5$ are $1$ modulo $4$. (b) Use the Chinese remainder theorem and Volume~2, Chapter~4 for the odd prime powers.
\end{hint}
\begin{hint}{ex:similarQ}
A Smith form of $XI-A$ computed in $K[X]$ is also one in $L[X]$.
\end{hint}
\begin{hint}{ex:nilpotent}
The invariant factors divide $\mu_N=X^k$. Compute $\dim\ker x^k$ on $K[X]/(X^a)$.
\end{hint}
\begin{hint}{ex:diagonalisable}
$K[X]/(X-\lambda)$ is one-dimensional, with $x$ acting as $\lambda$. For $A^k=I$, $X^k-1$ has distinct roots in $\C$.
\end{hint}
\begin{hint}{ex:cyclicvector}
(i) means $\deg d_r=n$, i.e.\ there is only one invariant factor. For (ii), what is the annihilator of $v$?
\end{hint}
\begin{hint}{ex:sublattices}
Imitate Lemma~\ref{lem:subfree} (or the proof of the analogous lemma in Volume~13, Chapter~1): look at the second coordinates, then at the elements with second coordinate $0$. The index is a determinant (Corollary~\ref{cor:index}).
\end{hint}

\begin{hint}{pb:kronecker}
(a) If $C$ is the companion matrix of $P$, the characteristic polynomial of $C^k$ has integer coefficients and roots
$\alpha_i^k$. There are only finitely many monic integer polynomials of degree $d$ with all roots in the closed unit disc.
(b) Look at $x^dP(x+\frac1x)$.
\end{hint}

\section{Solutions}

\begin{solution}{ex:ab360}
\idea Classify each Sylow subgroup separately (one partition per prime), then recombine.
\writeup The $2$-part has order $8$: $\Z/8$, $\Z/4\oplus\Z/2$ or $(\Z/2)^3$; the $3$-part has order $9$: $\Z/9$ or $(\Z/3)^2$; the $5$-part is $\Z/5$. The six combinations, with their invariant factors (multiply the largest powers of each prime, then the next largest, \dots):
\begin{center}
\small
\begin{tabular}{ll}
elementary divisors & invariant factors\\\hline
$8,9,5$ & $\Z/360$\\
$8,3,3,5$ & $\Z/3\oplus\Z/120$\\
$4,2,9,5$ & $\Z/2\oplus\Z/180$\\
$4,2,3,3,5$ & $\Z/6\oplus\Z/60$\\
$2,2,2,9,5$ & $\Z/2\oplus\Z/2\oplus\Z/90$\\
$2,2,2,3,3,5$ & $\Z/2\oplus\Z/6\oplus\Z/30$
\end{tabular}
\end{center}
\end{solution}

\begin{solution}{ex:smithcalc}
\idea The invariant factors can be read from gcds of minors, without performing the operations.
\writeup $\Delta_1=\gcd(4,6,6,4)=2$ and $\Delta_2=|\det A|=|16-36|=20$. So $d_1=2$, $d_2=20/2=10$, and by Corollary~\ref{cor:index} the group is $\Z/2\oplus\Z/10$, of order $20$. Check with operations: $R_2\leftarrow R_2-R_1$ gives $\begin{psmallmatrix}4&6\\2&-2\end{psmallmatrix}$; swap the rows and clear: $\begin{psmallmatrix}2&-2\\4&6\end{psmallmatrix}\to\begin{psmallmatrix}2&-2\\0&10\end{psmallmatrix}\to\begin{psmallmatrix}2&0\\0&10\end{psmallmatrix}$.
\end{solution}

\begin{solution}{ex:unitsmodn}
\idea For odd prime powers the unit group is cyclic (Volume~2); for $2^k$, track the $2$-adic size of $5^{2^j}-1$, as in Volume~2's lemma for odd $p$.
\writeup (a) $5=1+2^2\cdot1$, so $t_0=1$. If $5^{2^j}=1+2^{j+2}t_j$ with $t_j$ odd, then $5^{2^{j+1}}=1+2^{j+3}t_j+2^{2j+4}t_j^2=1+2^{j+3}(t_j+2^{j+1}t_j^2)$, and $t_{j+1}=t_j+2^{j+1}t_j^2$ is odd. So $5^{2^{k-3}}\not\equiv1$ and $5^{2^{k-2}}\equiv1\pmod{2^k}$: the order of $5$, a power of $2$ dividing $\varphi(2^k)=2^{k-1}$, is $2^{k-2}$. Powers of $5$ are $\equiv1\pmod4$, so $-1\notin\langle5\rangle$, and $\langle-1\rangle\cap\langle5\rangle=\{1\}$. The group is abelian, so $\langle-1\rangle\langle5\rangle$ is a subgroup isomorphic to $\Z/2\times\Z/2^{k-2}$ (Proposition~\ref{prop:internal}), of order $2^{k-1}=|(\Z/2^k\Z)^*|$: it is everything.

(b) $360=8\cdot9\cdot5$, so $(\Z/360\Z)^*\cong(\Z/8)^*\times(\Z/9)^*\times(\Z/5)^*\cong(\Z/2)^2\times\Z/6\times\Z/4$ (Theorem~\ref{thm:crtring}, (a), and Volume~2, Chapter~4 for $9$ and $5$). Elementary divisors: $2,2,2,4,3$. Invariant factors: $d_4=4\cdot3=12$, $d_3=d_2=d_1=2$: $(\Z/360\Z)^*\cong\Z/2\oplus\Z/2\oplus\Z/2\oplus\Z/12$, of order $96=\varphi(360)$. In $\bigoplus\Z/d_i$ the smallest $e$ with $ex=0$ for all $x$ is $d_r$, so $e=12$: $a^{12}\equiv1\pmod{360}$ for every $a$ prime to $360$, while Euler's theorem only gives $a^{96}\equiv1$. (This $e$ is Carmichael's function $\lambda(360)$.)
\end{solution}

\begin{solution}{ex:similarQ}
\idea Invariant factors are computed by operations inside $K[X]$, and are unique, so they cannot change when the field grows.
\writeup Let
\[
P(XI-A)Q=\operatorname{diag}(1,\dots,1,d_1,\dots,d_r)
\]
with $P,Q\in GL_n(K[X])$ and $d_i\in K[X]$ monic, $d_1\mid\dots\mid d_r$. The same $P,Q$ lie in $GL_n(L[X])$ (their determinants are non-zero constants), and divisibility in $K[X]$ implies divisibility in $L[X]$; so this is also a Smith form of $XI-A$ over $L[X]$, and by uniqueness the invariant factors of $A$ over $L$ are the $d_i$. If $A$ and $B$ are similar over $L$, they have the same invariant factors over $L$ (Theorem~\ref{thm:frobenius}), hence over $K$, hence they are similar over $K$.
\end{solution}

\begin{solution}{ex:nilpotent}
\idea Nilpotent means all invariant factors are powers of $X$; the exponents form a partition.
\writeup If $N^k=0$, $\mu_N$ divides $X^k$, so every invariant factor is $X^{a_i}$, with $a_1\le\dots\le a_r$ and $\sum a_i=\deg\chi_N=n$: a partition of $n$. By Theorem~\ref{thm:frobenius}, $N$ is determined up to similarity by this partition, and every partition occurs (take $\operatorname{diag}(C(X^{a_1}),\dots)$). On $K[X]/(X^a)$, $\ker x^k$ consists of the classes of multiples of $X^{\max(a-k,0)}$, of dimension $\min(a,k)$. So $\dim\ker N^k=\sum_i\min(a_i,k)$, and $\dim\ker N^k-\dim\ker N^{k-1}=\#\{i:a_i\ge k\}$: these numbers determine the partition. For $n=4$: the partitions $4$, $3+1$, $2+2$, $2+1+1$, $1+1+1+1$ give five classes, with $\dim\ker N=1,2,2,3,4$ and the zero matrix last.
\end{solution}

\begin{solution}{ex:diagonalisable}
\idea Diagonalisable means $V_A$ is a sum of one-dimensional modules $K[X]/(X-\lambda)$; the invariant factors then all divide a product of distinct linear factors.
\writeup If
\[
A=S\operatorname{diag}(\lambda_1,\dots,\lambda_n)S^{-1},
\]
then $\prod_{\lambda\text{ distinct}}(X-\lambda)$ kills $A$, and any polynomial killing $A$ vanishes at every $\lambda_i$; so $\mu_A$ is this product. Conversely, if $\mu_A=d_r$ is a product of distinct $X-\lambda$, so is every $d_i$ (it divides $d_r$), and by the Chinese remainder theorem $K[X]/(d_i)\cong\bigoplus K[X]/(X-\lambda)$, each summand being a line on which $A$ acts as $\lambda$: $V_A$ has a basis of eigenvectors. If $A^2=A$, $\mu_A$ divides $X^2-X=X(X-1)$. If $A^k=I$ over $\C$, $\mu_A$ divides $X^k-1=\prod_{j}(X-e^{2i\pi j/k})$, with distinct roots.
\end{solution}

\begin{solution}{ex:cyclicvector}
\idea A module is cyclic (generated by one element) iff it has a single invariant factor.
\writeup (i)$\Leftrightarrow$(iii): $\chi_A=d_1\cdots d_r$ and $\mu_A=d_r$, so (i) means that there is a single invariant factor $d_1=\chi_A$, i.e.\ $A$ is similar to $C(\chi_A)$ (Theorem~\ref{thm:frobenius}); and a companion matrix $C(d)$ has the single invariant factor $d$, since $V_{C(d)}\cong K[X]/(d)$. (iii)$\Rightarrow$(ii): in $K[X]/(d)$, $v=1$ works. (ii)$\Rightarrow$(i): the map $K[X]\to V_A$, $g\mapsto g(A)v$, is onto, with kernel $(m)$ for some monic $m$; so $V_A\cong K[X]/(m)$, which has a single invariant factor, and (i) holds by the first equivalence.
\end{solution}

\begin{solution}{ex:sublattices}
\idea A subgroup of $\Z^2$ has a triangular basis (Hermite normal form), which can be normalised uniquely.
\writeup Let $L$ have index $n$. Its second coordinates form a subgroup $c\Z$ of $\Z$, and $c\ge1$: otherwise $L\subseteq\Z\times\{0\}$ and $\Z^2/L$ would be infinite. Its elements with second coordinate $0$ form $a\Z\times\{0\}$, with $a\ge1$ since $(n,0)\in L$ ($n$ kills $\Z^2/L$ by Lagrange). Pick $(b,c)\in L$; as in Lemma~\ref{lem:subfree}, $L=\Z(a,0)+\Z(b,c)$, and $b$ may be changed by multiples of $a$, so we can take $0\le b<a$. Then $a$ and $c$ are determined by $L$, and so is $b$ in $[0,a)$: if $(b,c),(b',c)\in L$ then $(b-b',0)\in L$, so $a\mid b-b'$. By Corollary~\ref{cor:index}, the index is $\left|\det\begin{psmallmatrix}a&b\\0&c\end{psmallmatrix}\right|=ac$. Conversely each triple with $ac=n$, $0\le b<a$ gives a subgroup of index $n$, and different triples give different subgroups. For each divisor $a$ of $n$ there are $a$ choices of $b$: in total $\sum_{a\mid n}a=\sigma(n)$. For $n=2$: three subgroups, the kernels of the three non-zero homomorphisms $\Z^2\to\Z/2$.
\end{solution}

\begin{solution}{pb:kronecker}
\idea Raising the roots to the $k$-th power keeps them in the closed unit disc and keeps the polynomial in $\Z[X]$, of the
same degree. But there are only finitely many such polynomials: their coefficients are bounded by Viète. So the powers
$\alpha^k$ take finitely many values, and $\alpha^k=\alpha^l$ for some $k<l$. To see that $\prod(X-\alpha_i^k)$ has integer
coefficients without symmetric functions, take a matrix.

\writeup (a) Let $d=\deg P$, with roots $\alpha_1,\dots,\alpha_d$ (with multiplicity), and $C\in M_d(\Z)$ its companion
matrix, whose characteristic polynomial is $P$. Over $\C$, $C$ is similar to an upper triangular matrix $T$ with diagonal
$\alpha_1,\dots,\alpha_d$ (Jordan form, Volume~3, Chapter~7). Then $C^k$ is similar to $T^k$, upper triangular with diagonal
$\alpha_i^k$, so $P_k=\prod_i(X-\alpha_i^k)$ is the characteristic polynomial of $C^k\in M_d(\Z)$: $P_k\in\Z[X]$, monic of degree
$d$. Its coefficients are, up to sign, the elementary symmetric functions of the $\alpha_i^k$ (Volume~1, Chapter~7), and since
$|\alpha_i^k|\le1$, the coefficient of $X^{d-j}$ has modulus at most $\binom dj$. There are finitely many monic integer
polynomials of degree $d$ with these bounds, so the set $R$ of all their roots is finite. For each $i$, the numbers
$\alpha_i,\alpha_i^2,\alpha_i^3,\dots$ all lie in $R$, so $\alpha_i^k=\alpha_i^l$ for some $k<l$; as $\alpha_i\ne0$ (because
$P(0)\ne0$), $\alpha_i^{l-k}=1$.

If $P$ is irreducible, its root $\alpha_1$ is a root of unity, of some order $n$, so it is a root of $\Phi_n\in\Z[X]$, which is
irreducible (Volume~3, Chapter~10). Two monic irreducible polynomials of $\Q[X]$ with a common root are equal: $P=\Phi_n$.

(b) Let $Q(x)=x^dP(x+\frac1x)$. Expanding $(x+\frac1x)^j=\sum_i\binom jix^{j-2i}$, $Q$ is a monic integer polynomial of
degree $2d$ with $Q(0)=1$. If $\beta=2\cos\theta$ is a root of $P$, then $x+\frac1x=\beta$ has the roots $e^{\pm i\theta}$, of
modulus $1$. Since $P=\prod(X-\beta_i)$, $Q=\prod_i(x^2-\beta_ix+1)$, and all its roots are of this form. By (a) they are
roots of unity: $e^{i\theta}=e^{2i\pi r}$ with $r\in\Q$, and $\beta=2\cos2\pi r$.

(c) $z+\bar z=\frac65$ and $z\bar z=1$, so $z$ is a root of $5X^2-6X+5$, irreducible over $\Q$ (negative discriminant).
If $z$ were a root of unity, it would be a root of some $\Phi_n$, so $5X^2-6X+5$ would divide $\Phi_n$ in $\Q[X]$; by Gauss's
lemma (Volume~1, Chapter~7) a primitive divisor of a monic integer polynomial is monic up to sign, and $5X^2-6X+5$ is
primitive and not monic. So $z$ is not a root of unity. What fails is integrality: the minimal polynomial is not in
$\Z[X]$ once made monic, $X^2-\frac65X+1$.
\tools Volume~1, Chapters~7 and~8 (Viète, Gauss's lemma, characteristic polynomials); Volume~3, Chapters~7 and~10 (Jordan form,
irreducibility of $\Phi_n$).
\leads Kronecker proved this in 1857. Lehmer asked in 1933 whether the \emph{Mahler measure} $\prod\max(1,|\alpha_i|)$ of a
non-cyclotomic irreducible integer polynomial (see Problem~\ref{pb:jensen}) is bounded away from $1$. The smallest known value,
$1.17628\ldots$, belongs to Lehmer's own polynomial $x^{10}+x^9-x^7-x^6-x^5-x^4-x^3+x+1$, and the question is still open
(Volume~13).
\end{solution}

\section*{Sources and further reading}

\begin{itemize}
\item H.~J.~S.~Smith, \emph{On systems of linear indeterminate equations and congruences}, Phil.\ Trans.\ R.\ Soc.\ London 151 (1861), 293--326. The normal form, for integer matrices.
\item G.~Frobenius, \emph{Theorie der linearen Formen mit ganzen Coefficienten}, J.\ reine angew.\ Math.\ 86 (1879), 146--208. Invariant factors and the classification of finite abelian groups.
\item D.~S.~Dummit and R.~M.~Foote, \emph{Abstract Algebra}, 3rd ed., Wiley, 2004, Chapter~12. Modules over a PID, the rational canonical form and the Jordan form.
\item R.~Kannan and A.~Bachem, \emph{Polynomial algorithms for computing the Smith and Hermite normal forms of an integer matrix}, SIAM J.\ Comput.\ 8 (1979), 499--507. Why the naive algorithm needs care: the entries can explode.
\item H.~Edelsbrunner and J.~Harer, \emph{Computational Topology: An Introduction}, AMS, 2010. Homology computed by matrix reduction.
\item A.~Cayley, \emph{A memoir on the theory of matrices}, Phil.\ Trans.\ R.\ Soc.\ London 148 (1858), 17--37, where the quotation appears.
\end{itemize}
